% !PRINTED 708
\weberpage{708}
One obtains the numbers $n$ if one removes from the totality of all natural
numbers $k$ the multiples of $q$, that is, the totality of the numbers $qk$.
Accordingly
\[
\Sigma \frac{1}{n^s} = \Sigma \frac{1}{k^s} - \Sigma \frac{1}{q^s k^s} = (1-q^{-s}) \Sigma \frac{1}{k^s}.
\]
If, however, in Theorem 3., \S.\,191, one puts for the $\mu_n$ the sequence of
the natural numbers $k$, it follows that:
\[
\underset{s=1}{\mathrm{Lim}} (s-1) \Sigma \frac{1}{k^s} = 1,
\]
and we therefore obtain
\weq{12}{\underset{s=1}{\mathrm{Lim}} \frac{s-1}{1-q^{-s}} \Sigma \frac{1}{n^s} = 1.}
The other factors of the product $\Phi(s)$, however, are continuous functions of
$s$ by Theorem 1., \S.\,191, if the infinite series are ordered according to
increasing values of $n$.

For the sum
\weq{13}{\Sigma \chi(n)}
ordered according to increasing values of $n$ is indeed not convergent, but it
cannot for all that exceed a finite limit in absolute value. For as often as $n$
runs through a complete system of residues with respect to the modulus $m$,
there is added to the sum (13) a contribution which by \S.\,11, 6. has the value
0. Accordingly, if $\chi$ is not the principal character,
\weq{14}{\underset{s=1}{\mathrm{Lim}} \overset{n}{\Sigma} \frac{\chi(n)}{n^s} = \overset{n}{\Sigma} \frac{\chi(n)}{n},}
and accordingly we now obtain for the class number $h$ in the cyclotomic field
$\Omega_m$ the expression
\weq{15}{g\,h = \overset{\chi}{\Pi} \overset{n}{\Sigma} \frac{\chi(n)}{n},}
in which the product sign $\Pi$ refers only to those characters $\chi$ which are
\gesperrt{different from the principal character}.

\begin{center}
\S.\,194.

\gesperrt{Determination of the Sums X.}
\end{center}

The infinite series
\weq{1}{X = \overset{n}{\Sigma} \frac{\chi(n)}{n},}
%CONT>>
